\section{Implemented restarted GMRES algorithm}

We solve
\begin{equation}
    A x = b, \qquad A\in\mathbb{C}^{n\times n},
    \quad x,b\in\mathbb{C}^{n},
\end{equation}
without preconditioning. Superscript $H$ denotes conjugate transpose, a bar denotes
scalar complex conjugation, and $\mathrm{i}^2=-1$. Vector norms are Euclidean:
$\lVert z\rVert_2=(\sum_{\ell=1}^n |z_\ell|^2)^{1/2}$. The implementation
accumulates norms with successive \texttt{hypot} calls rather than squaring all
entries directly. At the start of each restart cycle, let $x_{\mathrm{base}}$
denote the current approximation and set $r=b-Ax_{\mathrm{base}}$ and
$\beta=\lVert r\rVert_2$; in the first cycle $x_{\mathrm{base}}=x_0$ is the
user's initial guess. The stopping test is applied before normalization;
only if another iteration is required is $v_1=r/\beta$ formed. Complex vectors
use the Hermitian inner product
\begin{equation}
    \langle v,w\rangle = v^{H}w
    = \sum_{i=1}^{n}\overline{v_i}\,w_i.
    \label{eq:hermitian-product}
\end{equation}
The conjugation of the first argument in Eq.~\eqref{eq:hermitian-product} is
explicit in both the production implementation and the independent test
diagnostics.

Within a cycle, $j$ is the one-based Arnoldi step and
$V_j=[v_1,\ldots,v_j]\in\mathbb{C}^{n\times j}$. At step $j$, the implementation
initializes $w=Av_j$ and $h_{ij}=0$ for $1\leq i\leq j$. For each of two complete
modified Gram--Schmidt (MGS) passes, and sequentially for $i=1,\ldots,j$, it uses
\begin{align}
    \alpha &\leftarrow \langle v_i,w\rangle,\\
    h_{ij} &\leftarrow h_{ij}+\alpha,\qquad
    w\leftarrow w-\alpha v_i.
\end{align}
The same saved projection $\alpha$ is used in both updates. Set
$h_{j+1,j}=\lVert w\rVert_2$. If no near-breakdown is detected, the new vector is
$v_{j+1}=w/h_{j+1,j}$. In exact arithmetic, away from breakdown,
$AV_j=V_{j+1}\underline{H}_j$, where
$\underline{H}_j\in\mathbb{C}^{(j+1)\times j}$ is upper Hessenberg. The underline
denotes the rectangular Hessenberg matrix, not conjugation. This relation is
subject to rounding errors in the implementation. Reorthogonalization is a safeguard in
this implementation; finite-precision properties of MGS-GMRES and loss of
orthogonality are treated in detail by Paige, Rozlo\v{z}n\'{i}k, and
Strako\v{s} \cite{paige_rozloznik_strakos_2006}.

The Arnoldi relation gives the reduced least-squares problem
\begin{equation}
    y_j = \arg\min_{y\in\mathbb{C}^j}
    \left\lVert \beta e_1-\underline{H}_j y\right\rVert_2,
    \qquad x_j=x_{\mathrm{base}}+V_j y_j.
    \label{eq:least-squares}
\end{equation}
Here $e_1\in\mathbb{R}^{j+1}$ is the first coordinate vector, and $x_j$ is the
candidate within the current cycle. The upper-Hessenberg matrix is reduced by
left multiplication with complex unitary Givens rotations
\begin{equation}
    G = \begin{bmatrix}c&s\\-s^{\ast}&c\end{bmatrix},
    \qquad c\in\mathbb{R},\quad s\in\mathbb{C},\quad c^2+|s|^2=1.
    \label{eq:givens}
\end{equation}
The scalar $s^{\ast}=\overline{s}$ is explicitly conjugated; thus $G^H G=I_2$.
For entries $a,b\in\mathbb{C}$ to be rotated, the code sets
$d=\operatorname{hypot}(|a|,|b|)$ and, when both entries are nonzero,
\begin{equation}
    c=\frac{|a|}{d},\qquad
    s=\frac{a}{|a|}\frac{\overline{b}}{d},\qquad
    G\begin{bmatrix}a\\b\end{bmatrix}
      =\begin{bmatrix}(a/|a|)d\\0\end{bmatrix}.
\end{equation}
If $b=0$, it uses $(c,s)=(1,0)$, including $a=b=0$. If $a=0$ and
$b\ne0$, the code returns
\[
    c=0,\qquad s=\frac{\overline{b}}{|b|},\qquad
    G\begin{bmatrix}0\\b\end{bmatrix}
      =\begin{bmatrix}sb\\0\end{bmatrix}
      =\begin{bmatrix}\overline{b}b/|b|\\0\end{bmatrix}
      =\begin{bmatrix}|b|\\0\end{bmatrix}.
\]
The conjugate in $s$ fixes the phase of the first entry to the nonnegative real
number $|b|$. For example, $b=\mathrm{i}$ gives $s=-\mathrm{i}$ and
$s b=(-\mathrm{i})\mathrm{i}=1$. In contrast, $s=b/|b|$ would produce
$b^2/|b|$ (equal to $-1$ for $b=\mathrm{i}$) and is not the convention used by
\texttt{make\_rotation} and \texttt{rotate}.
Previous rotations are applied to each new Hessenberg column, and the new
rotation is applied both to that column and to the reduced right-hand side.
Triangular back substitution then solves the nonsingular reduced system.

A restart cycle allows at most $\min(m,n,N_{\max}-N)$ Arnoldi steps, where $m$
is the requested restart length, $N$ is the total completed iteration count at
cycle entry, and $N_{\max}$ is the global limit. Convergence or breakdown may
end a cycle earlier. The last approximation of a completed cycle becomes the
next cycle's base point.

\subsection{Stopping and diagnostic rules}

Use a common normalization for recomputed and estimated residuals:
\begin{equation}
    D_b=\begin{cases}\lVert b\rVert_2,&b\ne0,\\1,&b=0,\end{cases}
    \qquad \rho(x)=\frac{\lVert b-Ax\rVert_2}{D_b}.
    \label{eq:true-residual}
\end{equation}
For $b\ne0$, this is relative; for the exact zero right-hand side it is absolute.
Let $g^{(j)}$ be the reduced vector $\beta e_1$ after the first $j$ embedded
Givens rotations. The least-squares residual estimate is
\begin{equation}
    \widehat{\rho}_j=
    \frac{|g^{(j)}_{j+1}|}{D_b}.
\end{equation}
The solver never uses this estimate to declare convergence. It recomputes
$b-Ax_j$ in double precision with a CSR product after each candidate update and
returns \texttt{CONVERGED} only if the computed $\rho(x_j)\leq\tau$, where
$\tau>0$ is the finite user tolerance. The same check at the initial guess can
return success with zero iterations. These computed residuals, including the
initial value, form the public residual history; they are not exact-arithmetic
residuals or the independent wider-precision test values.

Near-breakdown is flagged when
\begin{equation}
    \lVert w\rVert_2 \leq
    16\,\epsilon_{\mathrm{double}}\,\lVert Av_j\rVert_2.
\end{equation}
Here $w$ is the remainder after both MGS passes, and
$\epsilon_{\mathrm{double}}=2^{-52}$ is the spacing above one in the reported
binary64 environment. If back substitution succeeds, convergence is checked
before this flag terminates the solve: a residual meeting the tolerance gives
happy breakdown; otherwise the status is \texttt{BREAKDOWN}. An exactly zero
diagonal in the rotated triangular problem also returns \texttt{BREAKDOWN},
retaining the previous approximation and repeating its residual in the history.
Exhausting $N_{\max}$ returns \texttt{NOT\_CONVERGED}. Invalid dimensions or
options are rejected, and detected nonfinite arithmetic raises an exception
rather than returning success.
